\section{המכניקה הקוונטית של גוף צפיד \LR{(The Quantum Top)}}

בשיעור זה נעסוק בתיאור הקוונטי של גוף צפיד (\LR{Rigid Body}). המוטיבציה הפיזיקלית היא תיאור מערכות מורכבות כמו מולקולות או גרעינים, אשר ניתן להתייחס אליהן בקירוב טוב כגוף שבו המרחקים בין החלקיקים קבועים, והדינמיקה העיקרית היא סיבוב במרחב.

\subsection{קינמטיקה קלאסית של גוף צפיד}

נניח מערכת המורכבת מאוסף חלקיקים. נגדיר שתי מערכות צירים:
\begin{enumerate}
    \item \textbf{מערכת המעבדה (\LR{Lab/Space Frame}):} מערכת אינרציאלית קבועה. הקואורדינטות יסומנו באינדקסים $i, j, k$ (למשל $x_i$).
    \item \textbf{מערכת הגוף (\LR{Body Frame}):} מערכת הצמודה לגוף ומסתובבת איתו. הקואורדינטות יסומנו באינדקסים $a, b, c$ (למשל $r_a$).
\end{enumerate}

בגוף צפיד, המיקום של חלקיק $n$ במערכת הגוף, $\vec{r}^{(n)}$, הוא קבוע. המיקום במערכת המעבדה, $\vec{x}^{(n)}(t)$, משתנה בזמן עקב הסיבוב. הקשר ביניהם נתון על ידי מטריצת סיבוב אורתוגונלית $R(t)$:

\begin{definitionbox}{טרנספורמציה בין מערכת גוף למעבדה}
המיקום של חלקיק $n$ במערכת המעבדה נתון על ידי:
\begin{equation}
    x^{(n)}_i(t) = \sum_a R_{ia}(t) r^{(n)}_a
\end{equation}
כאשר $R(t)$ היא מטריצת סיבוב המקיימת $R R^T = \mathbb{I}$.
\end{definitionbox}

\subsubsection{אנרגיה קינטית וטנזור האינרציה}
האנרגיה הקינטית הכוללת היא סכום האנרגיות של כל החלקיקים:
\begin{equation}
    T = \frac{1}{2} \sum_n m_n \left( \dot{\vec{x}}^{(n)} \right)^2 = \frac{1}{2} \sum_n m_n \dot{x}^{(n)}_i \dot{x}^{(n)}_i
\end{equation}
נציב את הנגזרת של המיקום ($\dot{x}_i = \dot{R}_{ia} r_a$):
\begin{equation}
    T = \frac{1}{2} \sum_n m_n (\dot{R}_{ia} r^{(n)}_a) (\dot{R}_{ib} r^{(n)}_b)
\end{equation}
(כאן השתמשנו בהסכם הסכימה של איינשטיין על אינדקסים כפולים).
נבצע סכימה על כל החלקיקים $n$ ונגדיר טנזור עזר סימטרי $N_{ab}$ התלוי רק בגיאומטריה של הגוף:
\begin{equation}
    N_{ab} = \sum_n m_n r^{(n)}_a r^{(n)}_b
\end{equation}
כעת הביטוי לאנרגיה הופך להיות:
\begin{equation}
    T = \frac{1}{2} \sum_{i,a,b} \dot{R}_{ia} \dot{R}_{ib} N_{ab} = \frac{1}{2} \text{Tr} \left( \dot{R} N \dot{R}^T \right)
\end{equation}

\begin{notebox}{מהירות זוויתית ומטריצות אנטי-סימטריות}
מכיוון ש-$R R^T = \mathbb{I}$, גזירה בזמן נותנת:
$$ \dot{R} R^T + R \dot{R}^T = 0 \implies \dot{R} R^T = - (\dot{R} R^T)^T $$
כלומר, המטריצה $\Omega \equiv \dot{R} R^T$ היא אנטי-סימטרית. במרחב תלת-ממדי, מטריצה אנטי-סימטרית מיוצגת על ידי 3 רכיבים (וקטור המהירות הזוויתית $\vec{\omega}$):
$$ \Omega_{ij} = \epsilon_{ijk} \omega_k $$
\end{notebox}

לאחר מניפולציות אלגבריות (שימוש בזהויות של טנזור לוי-צ'יביטה $\epsilon_{ijk}$), ניתן לכתוב את האנרגיה הקינטית באמצעות טנזור האינרציה המוכר $I_{ab}$:
\begin{equation}
    T = \frac{1}{2} \sum_{a,b} I_{ab} \omega_a \omega_b
\end{equation}
כאשר $\omega_a$ הם רכיבי המהירות הזוויתית במערכת הגוף, ו-$I_{ab}$ הוא טנזור האינרציה:
\begin{equation}
    I_{ab} = \sum_n m_n \left( (r^{(n)})^2 \delta_{ab} - r^{(n)}_a r^{(n)}_b \right)
\end{equation}

\subsection{תנע זוויתי: מעבדה מול גוף}

זהו החלק הקריטי להבנת המערכת הקוונטית. אנו מגדירים שני וקטורי תנע זוויתי:

\begin{enumerate}
    \item \textbf{במערכת המעבדה ($L_i$):} זהו התנע הזוויתי הפיזיקלי הנמדד על ידי צופה חיצוני.
    \begin{equation}
        L_i = \sum_n m_n (\vec{x}^{(n)} \times \dot{\vec{x}}^{(n)})_i
    \end{equation}
    \item \textbf{במערכת הגוף ($K_a$):} זהו התנע הזוויתי כפי שהוא מוטל על צירי הגוף המסתובבים.
    \begin{equation}
        K_a = \sum_i R_{ia} L_i \quad \text{(הטלה של } \vec{L} \text{ על צירי הגוף)}
    \end{equation}
\end{enumerate}

באמצעות פיתוח מתמטי (המופיע בפירוט בהרצאה), ניתן לקשר את התנעים הללו למהירות הזוויתית ולמומנט האינרציה. הקשר החשוב ביותר הוא ההמילטוניאן:

\begin{resultbox}{ההמילטוניאן של גוף צפיד}
מכיוון שהאנרגיה היא קינטית בלבד, ההמילטוניאן ניתן לכתיבה באמצעות אופרטורי התנע הזוויתי במערכת הגוף ($K$) וטנזור האינרציה ההופכי:
\begin{equation}
    H = \frac{1}{2} \sum_{a,b} (I^{-1})_{ab} K_a K_b
\end{equation}
\end{resultbox}

\subsection{קוונטיזציה ויחסי חילוף}

כדי לעבור למכניקת הקוונטים, אנו מתייחסים ל-$R_{ia}$ כמשתנים הדינמיים (הקואורדינטות המוכללות). הקוונטיזציה מתבצעת על ידי הגדרת יחסי חילוף קנוניים. התוצאות המרכזיות הן יחסי החילוף של אופרטורי התנע הזוויתי בשתי המערכות.

\subsubsection{יחסי חילוף במערכת המעבדה ($L$)}
אופרטורי $L_i$ הם יוצרי הסיבובים במרחב המעבדה. הם מקיימים את האלגברה המוכרת של $SO(3)$:
\begin{equation}
    [L_i, L_j] = i\hbar \epsilon_{ijk} L_k
\end{equation}

\subsubsection{יחסי חילוף במערכת הגוף ($K$)}
כאן מתקבלת תוצאה מפתיעה וחשובה. לאחר חישוב יחסי החילוף בין הרכיבים המוטלים על מערכת הגוף, מקבלים סימן מינוס "חריג":

\begin{resultbox}{יחסי חילוף אנומליים של $K$}
\begin{equation}
    [K_a, K_b] = -i\hbar \epsilon_{abc} K_c
\end{equation}
\textbf{הסבר:} הסימן ההפוך נובע מכך שסיבוב המערכת "גוף" יחסית למעבדה הוא הפוך לסיבוב המערכת "מעבדה" יחסית לגוף. מתמטית, זה נובע מסדר האינדקסים בטרנספורמציה $R$.
\end{resultbox}

\subsubsection{יחסים בין המערכות}
שתי מערכות התנע הזוויתי מתחלפות זו עם זו:
\begin{equation}
    [L_i, K_a] = 0
\end{equation}
כמו כן, ריבוע התנע הזוויתי הוא אינווריאנטי (סקלר), ולכן זהה בשתי המערכות:
\begin{equation}
    L^2 = \sum_i L_i^2 = \sum_a K_a^2 = K^2
\end{equation}

\subsection{ספקטרום האנרגיה: סביבון סימטרי}

כדי למצוא את רמות האנרגיה, נבחר מערכת צירים בגוף המתלכדת עם הצירים הראשיים של טנזור האינרציה, כך ש-$I_{ab}$ אלכסוני עם ערכים עצמיים $I_x, I_y, I_z$. ההמילטוניאן הוא:
\begin{equation}
    H = \frac{K_x^2}{2I_x} + \frac{K_y^2}{2I_y} + \frac{K_z^2}{2I_z}
\end{equation}

נבחר בסיס של מצבים עצמיים משותפים לאופרטורים שמתחלפים עם $H$. האופרטורים הללו הם:
\begin{itemize}
    \item $L^2 = K^2$ (ערך עצמי $\hbar^2 L(L+1)$).
    \item $L_z$ (הטלה על ציר $z$ במעבדה, ערך עצמי $\hbar M$).
    \item $K_z$ (הטלה על ציר $z$ בגוף, ערך עצמי $\hbar k$).
\end{itemize}
נסמן את המצבים כ-$|L, M, k\rangle$.

\begin{examplebox}{מקרה פרטי: סביבון סימטרי (\LR{Symmetric Top})}
נניח סימטריה גלילית בגוף, כך ש-$I_x = I_y \neq I_z$. ההמילטוניאן הופך ל:
\begin{equation}
    H = \frac{K_x^2 + K_y^2}{2I_x} + \frac{K_z^2}{2I_z} = \frac{K^2 - K_z^2}{2I_x} + \frac{K_z^2}{2I_z}
\end{equation}
מכיוון ש-$K^2$ ו-$K_z$ אלכסוניים בבסיס שבחרנו, האנרגיות הן:
\begin{equation}
    E_{L,k} = \frac{\hbar^2}{2I_x} L(L+1) + \frac{\hbar^2}{2} \left( \frac{1}{I_z} - \frac{1}{I_x} \right) k^2
\end{equation}
\end{examplebox}

\begin{theorembox}{תובנה פיזיקלית: סקאלות אנרגיה}
במערכות כמו מולקולות או גרעינים, מומנט האינרציה $I$ הוא בדרך כלל גדול מאוד (יחסית למסת האלקטרון), שכן הוא תלוי במסת הגרעינים ($M_{nuc}$) ובמרחק ביניהם.
מכיוון שהאנרגיה הולכת כמו $1/I$, רמות האנרגיה הסיבוביות (\LR{Rotational Levels}) הן נמוכות מאוד בהשוואה לרמות האלקטרוניות או הוויברציוניות. זו הסיבה שהספקטרום הסיבובי שולט באנרגיות הנמוכות ומאפשר מדידה מדויקת של מבנה המולקולה/הגרעין דרך קבועי האינרציה.
\end{theorembox}